\chapter{A short source guide}
\label{ch:sources}

This book emphasizes explanations and calculations. The following
guide is deliberately selective: it gives a route to the main
constructions and to the strongest external inputs without turning
each page into a literature survey. The repository manuscripts
and the public literature have different roles. In particular,
the detailed lattice realization developed in Part II is presented
as the repository's research construction with the finite inputs
and qualifications stated in the text.

\raggedright

For the elementary theta and eta transformations, the
\href{https://dlmf.nist.gov/20.7#viii}{NIST theta-function formulas}
and \href{https://dlmf.nist.gov/23.18}{eta modular transformations}
fix the analytic conventions used in the rank-one lattice example.
The discriminant and cyclotomic ramification facts used in
Example~\ref{ex:coset-ramification} are explained in
J.~S.~Milne's \href{https://www.jmilne.org/math/CourseNotes/ANT.pdf}
{\emph{Algebraic Number Theory}}, Sections~3 and~6.

\section*{Fusion categories, centers, and modular data}
\addcontentsline{toc}{section}{Fusion categories, centers, and modular data}

D.~E.~Evans and T.~Gannon,
\emph{The exoticness and realisability of twisted Haagerup--Izumi
modular data},
\href{https://arxiv.org/abs/1006.1326}{arXiv:1006.1326},
is the principal source for the modular data and the
vector-valued character construction used here.
It is especially useful after Chapters~\ref{ch:modular}
and~\ref{ch:characters}, when the difference between modular
data and a realized VOA is clear.

M.~Izumi,
\emph{The classification of $3^n$ subfactors and related fusion
categories},
\href{https://arxiv.org/abs/1609.07604}{arXiv:1609.07604},
develops the operator-algebraic reconstruction and classification
framework. P.~Grossman and N.~Snyder,
\emph{Quantum subgroups of the Haagerup fusion categories},
\href{https://arxiv.org/abs/1102.2631}{arXiv:1102.2631},
is a source for the relations among the Haagerup categories and
their module categories.

For the center and extension mechanisms, useful primary sources are
M.~M\"uger,
\emph{From subfactors to categories and topology II},
\href{https://arxiv.org/abs/math/0111205}{arXiv:math/0111205},
and A.~Davydov, M.~M\"uger, D.~Nikshych, and V.~Ostrik,
\emph{The Witt group of non-degenerate braided fusion categories},
\href{https://arxiv.org/abs/1009.2117}{arXiv:1009.2117}.
The small pointed-center and algebra-object examples in this book
are intended to make those structural results easier to enter.

The cyclic-rotation definition and center-trace formula for
Frobenius--Schur indicators used in Chapter~\ref{ch:branching}
are due to S.-H.~Ng and P.~Schauenburg; see
\href{https://arxiv.org/abs/math/0601012}{\emph{Frobenius--Schur
indicators and exponents of spherical categories}}, Theorem~4.1.
The group and pointed-center calculations in that chapter fix
the normalization needed for the Haagerup twist moments.

D.~E.~Evans and T.~Gannon,
\emph{Non-unitary fusion categories and their doubles via
endomorphisms},
\href{https://arxiv.org/abs/1506.03546}{arXiv:1506.03546},
provides reconstruction machinery relevant to the arithmetic
chapter. The uniform arithmetic existence argument discussed there
comes from the repository research manuscript; it should not be
attributed wholesale to this paper.

For Chapter~\ref{ch:related-categories}, Grossman--Izumi--Snyder,
\href{https://arxiv.org/abs/1501.07324}{\emph{The Asaeda--Haagerup
fusion categories}}, supplies the published order-eight generalized
Haagerup example. Its exact hyperbolic-gamma identification is a local
result recorded in \sourcefile{hi_noncyclic_gis_gamma_proof.tex}.
The distinction between the completed order-twelve algebraic
construction and its open gamma identification is recorded in
\sourcefile{hi_noncyclic_gamma_identification_status.md}, with the
algebraic certificate explained in
\sourcefile{hi_noncyclic_ansatz_z2xz6_exact.tex}.

For extended Haagerup, the primary construction is
Bigelow--Morrison--Peters--Snyder,
\href{https://arxiv.org/abs/0909.4099}{\emph{Constructing the extended
Haagerup planar algebra}}. Grossman--Morrison--Penneys--Peters--Snyder,
\href{https://arxiv.org/abs/1810.06076}{\emph{The Extended Haagerup
fusion categories}}, develops its fusion-category and Morita theory;
Morrison--Snyder,
\href{https://arxiv.org/abs/1002.0168}{\emph{Non-cyclotomic fusion
categories}}, explains the field-of-definition distinction. The local
\sourcefile{extended_haagerup_reconstruction_and_coefficients.tex}
reorganizes the published coefficient seeds and explains their
reconstruction into categorical data. These category constructions
do not by themselves supply a new conformal field theory.

\section*{Rationality and the analytic bridge}
\addcontentsline{toc}{section}{Rationality and the analytic bridge}

T.~Abe, G.~Buhl, and C.~Dong,
\emph{Rationality, regularity, and $C_2$-cofiniteness},
\href{https://arxiv.org/abs/math/0204021}{arXiv:math/0204021},
is a basic source for the finiteness implications.
R.~McRae,
\emph{On rationality for $C_2$-cofinite vertex operator algebras},
\href{https://arxiv.org/abs/2108.01898}{arXiv:2108.01898},
and P.~Etingof and D.~Penneys,
\emph{Rigidity of non-negligible objects of moderate growth in
braided categories},
\href{https://arxiv.org/abs/2412.17681v3}{arXiv:2412.17681v3},
are the deeper inputs in the regularity route.
Their hypotheses are part of the argument; quoting their
conclusions without checking the relevant module category would
skip the main work of Chapter~\ref{ch:regularity}.

T.~Arakawa, C.~H.~Lam, and H.~Yamada,
\emph{Parafermion vertex operator algebras and $W$-algebras},
\href{https://arxiv.org/abs/1701.06229}{arXiv:1701.06229},
is useful for the parafermion identification in
Chapter~\ref{ch:factors}.
C.~Dong and Q.~Wang,
\emph{Quantum dimensions and fusion rules for parafermion vertex
operator algebras},
\href{https://arxiv.org/abs/1411.7429}{arXiv:1411.7429},
supplies the affine and parafermion branching, identification, and
fusion formulas used in that chapter's Potts-module example.

S.~Carpi, Y.~Kawahigashi, R.~Longo, and M.~Weiner,
\emph{From vertex operator algebras to conformal nets and back},
\href{https://arxiv.org/abs/1503.01260}{arXiv:1503.01260},
provides the locality and net framework.
A.~G.~Henriques and J.~E.~Tener,
\emph{Every conformal net has an associated unitary VOA},
\href{https://arxiv.org/abs/2507.20735v3}{arXiv:2507.20735v3},
supplies the reconstruction input in the version used by the
research manuscript. J.~E.~Tener,
\emph{Fusion and positivity in chiral conformal field theory},
\href{https://arxiv.org/abs/1910.08257}{arXiv:1910.08257},
is the source for the final fusion-positivity implication.
Chapter~\ref{ch:unitarity} separates the module-counting,
domain, and annulus-comparison steps needed to apply these results.

M.~Krauel and M.~Miyamoto,
\href{https://arxiv.org/abs/1501.00124}{\emph{A modular invariance
property of multivariable trace functions for regular vertex operator
algebras}}, Theorem~1, supplies the simultaneous two-temperature
transformation used in Chapter~\ref{ch:cosets}. The regular ambient
algebra and trace analyticity are the relevant hypotheses there;
the two smaller Virasoro algebras need not be rational.

B.~Gui,
\emph{Q-systems and extensions of completely unitary vertex
operator algebras},
\href{https://arxiv.org/abs/1908.03484}{arXiv:1908.03484},
and Y.-Z.~Huang, A.~Kirillov, Jr., and J.~Lepowsky,
\emph{Braided tensor categories and extensions of vertex operator
algebras},
\href{https://arxiv.org/abs/1406.3420}{arXiv:1406.3420},
explain the relation between categorical algebra objects and
actual extensions. These are the relevant sources when moving
from the center equivalence to the selected holomorphic theory.

For the actual diagonal full conformal net, K.-H.~Rehren,
\href{https://arxiv.org/abs/math/9911149}{\emph{Canonical tensor
product subfactors}}, gives the Longo--Rehren route used here
(Theorem~1.4, Proposition~1.5, and Corollary~1.6).
For categorical sewing, the paper by J.~Fjelstad, J.~Fuchs,
I.~Runkel, and C.~Schweigert,
\href{https://arxiv.org/abs/hep-th/0503194}{\emph{TFT construction
of RCFT correlators V}}, proves mapping-class invariance and
bulk and boundary factorization.
The paper by M.~Bischoff, Y.~Kawahigashi, and R.~Longo,
\href{https://arxiv.org/abs/1410.8848}{\emph{Characterization of
2D rational local conformal nets and its boundary conditions}},
identifies the categorical full center with the corresponding
net construction. Chapter~\ref{ch:sewing} distinguishes these
theorems from the analytic sewing statements used to interpret
surface amplitudes.

\section*{The physical chain and the central-charge-eight proposal}
\addcontentsline{toc}{section}{The physical chain and the central-charge-eight proposal}

For the controlled marginal directions,
S.~F\"orste and D.~Roggenkamp,
\emph{Current-current deformations of conformal field theories,
and WZW models},
\href{https://arxiv.org/abs/hep-th/0304234}{arXiv:hep-th/0304234},
gives the abelian current-deformation framework used in
Chapter~\ref{ch:full_theories}.

T.-C.~Huang, Y.-H.~Lin, K.~Ohmori, Y.~Tachikawa, and M.~Tezuka,
\emph{Numerical evidence for a Haagerup conformal field theory},
\href{https://arxiv.org/abs/2110.03008}{arXiv:2110.03008},
and R.~Vanhove and collaborators,
\emph{A critical lattice model for a Haagerup conformal field theory},
\href{https://arxiv.org/abs/2110.03532}{arXiv:2110.03532},
are two complementary starting points for the physical
motivation. Their lattice constructions and evidentiary roles
should be kept distinct.

The overlap approach developed in Part IV is related to
Y.~Liu, Y.~Zou, and S.~Ryu,
\emph{Operator fusion from wavefunction overlaps: Universal
finite-size corrections and application to Haagerup model},
\href{https://arxiv.org/abs/2203.14992}{arXiv:2203.14992}.
The book derives the normalization and finite-size tests it
uses, so the reader need not take an overlap value on trust
as a direct measurement of central charge.

P.~Calabrese and J.~Cardy,
\emph{Entanglement Entropy and Quantum Field Theory},
\href{https://arxiv.org/abs/hep-th/0405152}{arXiv:hep-th/0405152},
supplies the interval and massive single-cut entropy framework.
Chapter~\ref{ch:tensor_networks} derives the replica weights and
separates those geometries from the finite-entanglement scaling
hypothesis used for variational states.

For the rational-candidate tests in Chapter~\ref{ch:identification},
Dong--Ren--Xu, \href{https://arxiv.org/abs/1507.03306}
{\emph{On Orbifold Theory}}, Theorem~4.4 and Lemma~4.7, supplies
the orbifold dimension formula under the stated hypotheses.
Etingof--Nikshych--Ostrik, \href{https://arxiv.org/abs/math/0203060}
{\emph{On fusion categories}}, Corollary~8.14 and Proposition~8.27,
supplies Morita invariance and the weak-integrality dimension bound.
For $W_3$, Beltaos--Gannon, \href{https://arxiv.org/abs/1004.1205}
{\emph{The $W_N$ minimal model classification}}, Theorem~3.3,
classifies modular-invariant matrices; Edie-Michell--Gannon,
\href{https://arxiv.org/abs/2408.02794}
{\emph{Type II quantum subgroups for quantum $\mathfrak{sl}_N$.
II: Classification}}, Theorem~1.1, Remark~1.2 and Section~5.1,
provides the affine module-category input. The chapter explains the
additional category and full-center transfers rather than identifying
matrix classification with OPE sewing. The height-thirteen comparison
is a separate finite exact repository input, and all these exclusions
retain their stated chiral-algebra-preservation premises.

J.~Albert, Y.~Honda, J.~Kaidi, and Y.~Zheng,
\emph{Haagerup symmetry in $(E_8)_1$?},
\href{https://arxiv.org/abs/2512.08225v2}{arXiv:2512.08225v2},
gives the public central-charge-eight proposal relevant to this
project. The question mark is part of its title. A proposal
based on modular and spectral structure, the repository's
explicit lattice construction, and an identification of a
microscopic chain are three different achievements.

\section*{The construction record used for this textbook}
\addcontentsline{toc}{section}{The construction record used for this textbook}

The main algebraic record is
\sourcefile{hi_cft_c8_proof.tex}, with the structural and analytic
arguments in
\sourcefile{hi_cft_c8_manuscript_structure.tex} and
\sourcefile{hi_cft_c8_manuscript_analytic.tex}.
These files give the finite lattice gates and the global
deductions on which Part II is based. The review note
\sourcefile{hi_cft_c8_explicitness_review_20260928.md}
and the expanded material in
\sourcefile{c8_expanded_20260929/}
help expose intermediate steps that are compressed in the
shorter manuscript.

The spectrum, coupling, correlator, and defect records are in
\sourcefile{c8_physics_20260927/}. They are the immediate
research inputs for the detailed observables in Part III.
This textbook rederives the small examples and explains their
normalizations; it does not substitute a few successful
low-weight calculations for a reconstruction of every bulk
nonvacuum OPE.

The chain discussion uses the completed research summaries and
their linked evidence, especially
\sourcefile{hi_cft_20261001_status.md} and
\sourcefile{hi_cft_20261001_theory_handoff.md}.
The completed October 5 continuation is recorded in
\sourcefile{hi_cft_20261005_status.md},
\sourcefile{hi_cft_20261005_chain_handoff.md}, and
\sourcefile{hi_cft_20261005_theory_handoff.md}.
Its qualified variational sequence and completed operator and
correlation comparisons are included in the later chain chapters.
Earlier calculations retain their original finite scope; a later
qualified tensor does not retroactively qualify an earlier endpoint.

Formal verification has a separate scope. The authoritative
description of the current conditional development is
\sourcefile{lean/HaagerupIzumi/C8/TRUST_BOUNDARY.md}.
Appendix~\ref{app:verification} explains how to read a
trust boundary and why a theorem proved from named finite
inputs is not the same as a machine verification of those
inputs. The local audit notes in this textbook's
\sourcefile{editorial/} directory record independent
mathematical reviews of the exposition itself.

These paths identify the source snapshot in the accompanying
workspace. The textbook source bundle is self-contained for
typesetting, but it does not duplicate the entire research
repository or all of its numerical checkpoints.
